% Generated draft: requires mathematical and source review.
\documentclass[11pt]{article}
\usepackage[margin=25mm]{geometry}
\usepackage{fontspec}
\setmainfont{Latin Modern Roman}
\usepackage{amsmath,amssymb,mathtools,unicode-math}
\setmathfont{Latin Modern Math}
\usepackage{enumitem,needspace,xurl,hyperref}
\setlength{\parindent}{0pt}
\setlength{\parskip}{5pt}
\setlength{\emergencystretch}{3em}
\newcommand{\N}{\mathbb{N}}
\newcommand{\Z}{\mathbb{Z}}
\newcommand{\Q}{\mathbb{Q}}
\newcommand{\R}{\mathbb{R}}
\newcommand{\C}{\mathbb{C}}
\newcommand{\F}{\mathbb{F}}
\newcommand{\A}{\mathbb{A}}
\newcommand{\PP}{\mathbb{P}}
\newcommand{\E}{\mathbb{E}}
\newcommand{\eps}{\varepsilon}
\newcommand{\dd}{\,\mathrm{d}}
\begin{document}
\section*{UnsolvedMath ID: 20000994}
\textbf{Source problem code:} catalogue entry
\subsection{Definitions and mathematical statement}
Matt Baker asks for a chip-firing interpretation of the full Tutte polynomial, beyond Merino's $T_G(1,y)$ theorem [25]. The final extracted sentence, ``We made to use some additional geometric structure'', is ungrammatical and is not silently replaced by verified text. The aid answers the broad question using auxiliary embedding/root choices, while proving an obstruction to one particular choice-free statistic.

For a finite connected loopless graph with sink $q$, let $\operatorname{Rec}(G,q)$ be its recurrent stable sandpiles and use the Merino level convention
\[
 \operatorname{lev}(c)=\sum_{v\neq q}c(v)-|E(G)|+\deg(q).
\]
The Tutte polynomial may be defined by
\[
 T_G(x,y)=\sum_{A\subseteq E(G)}(x-1)^{r(E)-r(A)}(y-1)^{|A|-r(A)},
 \quad r(A)=|V(G)|-\#\operatorname{components}(V(G),A).
\]
The aid invokes Bernardi's rooted-embedding activity expansion and a level-preserving tree-to-recurrent-sandpile bijection $\eta$. With transported internal activity $J_\eta(c)$, it reports
\[
 T_G(x,y)=\sum_{c\in\operatorname{Rec}(G,q)}x^{J_\eta(c)}y^{\operatorname{lev}(c)}.
\]
The statistic depends on the declared auxiliary choices. For plane graphs, the summary describes the $x$-exponent as the level of the sandpile associated to the complementary dual spanning tree, while $y$ records primal level, under compatible dual choices and loop conventions when necessary. These data are not supplied in full detail.

The canonicity obstruction concerns a cycle $C_n$, $n\geq3$, rooted at a vertex. Its recurrent configurations are the all-one nonsink configuration of level one and the $n-1$ configurations with one zero and all other nonsink entries one, of level zero. The aid proves that no rooted-automorphism-invariant nonnegative-integer statistic $J$ can satisfy
\[
 T_{C_n}(x,y)=\sum_c x^{J(c)}y^{\operatorname{lev}(c)}
\]
with one positive monomial per configuration. Rooted reflection exchanges some zero-position configurations, forcing repeated monomials, whereas
\[
 T_{C_n}(x,y)=x+x^2+\cdots+x^{n-1}+y.
\]
This excludes that invariant positive-monomial form, not every possible chip-firing interpretation or weighted/cancelling expression.

An orientation of the rooted cycle repairs it explicitly. Number the nonsink vertices $1,\ldots,n-1$ along that orientation; assign $J=j$ to the configuration whose zero is at vertex $j$, and $J=0$ to the all-one configuration. The displayed full polynomial follows.

For a connected simple unicyclic graph with sink on its unique cycle and $b$ attached bridges, the aid reports the extension through the natural bridge-reduction correspondence with cycle recurrent configurations. The level is preserved and the exponent becomes $b+J_{\mathrm{cycle}}$. Accordingly
\[
 T_G(x,y)=x^b\bigl(x+x^2+\cdots+x^{n-1}+y\bigr),
\]
where $n$ is the cycle length. The zero-position description refers to the reduced cycle configuration, not literal zeros on attached trees. This explicit family and the existing choice-dependent theory answer different aspects of the original broad request.
\subsection{Short English statement}
The full Tutte polynomial already has a chip-firing expansion after auxiliary embedding choices. The aid shows why a particular automorphism-invariant monomial statistic cannot work even on rooted cycles. Choosing a cycle orientation supplies an explicit statistic, which extends over attached bridges.
\subsection{Sources}
\begin{itemize}
\item Local UnsolvedMath registry, permanent ID $20000994$, source code \texttt{catalogue entry}, supplied complete source record and research summary.
\item AIM workshop \emph{Generalizations of chip-firing and the critical group}, Problem (29), Matt Baker; supplied source \url{https://aimath.org/pastworkshops/chipfiringproblems.pdf}. Canonical local record \texttt{aim-combinatorics-notes.json}, \texttt{notes[118]}. No PDF was retrieved.
\item Merino's theorem [25], as cited in the original text; Bernardi's rooted-embedding activity expansion and level-preserving bijection, as cited by the aid. No titles, years or full bibliographic entries are supplied; the original literature field is empty.
\item Research attempt 1: partial result; broad problem solved at the recorded run, registry status solved and classification \texttt{SOLVED-IN-LITERATURE}; supplied rigor label ``proved'', novelty confidence ``low''. The machine-generated aid is not peer reviewed and requires independent verification before citation.
\end{itemize}
\end{document}
